\noindent
\begin{minipage}[t]{0.26\textwidth}
    % Cột lề trái: căn phải các nhãn ghi chú quy tắc để dóng thẳng hàng với khung ở cột chính
    \raggedleft\sffamily\bfseries\color{stewartgreen}
    \vspace*{270pt}
    Quy tắc lũy thừa\par
    \vspace{60pt}
    Quy tắc căn thức\par
\end{minipage}%
\hfill
\begin{minipage}[t]{0.71\textwidth}
    % Cột nội dung chính
    \textbf{(b)}\, Ta thấy rằng $\lim_{x \to 1} f(x) = 2$. Nhưng $\lim_{x \to 1} g(x)$ không tồn tại vì các giới hạn trái và phải khác nhau:
    \[
    \lim_{x \to 1^-} g(x) = -2 \qquad\qquad \lim_{x \to 1^+} g(x) = -1
    \]
    Do đó chúng ta không thể dùng Quy tắc 4 cho giới hạn cần tìm. Nhưng chúng ta có thể dùng Quy tắc 4 cho các giới hạn một phía:
    \begin{align*}
    \lim_{x \to 1^-} [f(x)g(x)] &= \lim_{x \to 1^-} f(x) \cdot \lim_{x \to 1^-} g(x) = 2 \cdot (-2) = -4\\[4pt]
    \lim_{x \to 1^+} [f(x)g(x)] &= \lim_{x \to 1^+} f(x) \cdot \lim_{x \to 1^+} g(x) = 2 \cdot (-1) = -2
    \end{align*}
    Các giới hạn trái và phải không bằng nhau, vì vậy $\lim_{x \to 1} [f(x)g(x)]$ không tồn tại.

    \textbf{(c)}\, Các đồ thị cho thấy rằng
    \[
    \lim_{x \to 2} f(x) \approx 1{,}4 \qquad\text{và}\qquad \lim_{x \to 2} g(x) = 0
    \]
    Bởi vì giới hạn của mẫu số bằng $0$, chúng ta không thể dùng Quy tắc 5. Giới hạn đã cho không tồn tại vì mẫu số tiến dần về $0$ trong khi tử số tiến dần về một số khác không.\hfill{\color{stewartcyan}\footnotesize$\blacksquare$}

    \vspace{3pt}
    Nếu chúng ta sử dụng liên tiếp Quy tắc tích với $g(x) = f(x)$, ta thu được quy tắc sau đây.

    % Khung 6: Quy tắc lũy thừa
    \begin{tcolorbox}[
        sharp corners,
        colback=white,
        colframe=stewartred,
        boxrule=0.75pt,
        left=10pt, right=10pt, top=5pt, bottom=5pt,
        before skip=6pt, after skip=7pt
    ]
    \textbf{\color{stewartred}6.}\ \, $\displaystyle \lim_{x \to a} [f(x)]^n = \left[ \lim_{x \to a} f(x) \right]^n$ \qquad\quad trong đó $n$ là một số nguyên dương
    \end{tcolorbox}

    Một tính chất tương tự, mà bạn được yêu cầu chứng minh trong Bài tập 2.5.69, nghiệm đúng đối với căn thức:

    % Khung 7: Quy tắc căn thức
    \begin{tcolorbox}[
        sharp corners,
        colback=white,
        colframe=stewartred,
        boxrule=0.75pt,
        left=10pt, right=10pt, top=5pt, bottom=6pt,
        before skip=6pt, after skip=7pt
    ]
    \textbf{\color{stewartred}7.}\ \, $\displaystyle \lim_{x \to a} \sqrt[n]{f(x)} = \sqrt[n]{\lim_{x \to a} f(x)}$ \qquad\quad trong đó $n$ là một số nguyên dương\\[5pt]
    \hspace*{16pt}[\text{Nếu } $n$ \text{ chẵn, ta giả thiết rằng } $\displaystyle\lim_{x \to a} f(x) > 0$.]
    \end{tcolorbox}

    Khi áp dụng bảy quy tắc giới hạn này, chúng ta cần sử dụng hai giới hạn đặc biệt:

    % Khung 8 & 9
    \begin{tcolorbox}[
        sharp corners,
        colback=white,
        colframe=stewartred,
        boxrule=0.75pt,
        left=10pt, right=10pt, top=6pt, bottom=6pt,
        before skip=6pt, after skip=7pt
    ]
    \textbf{\color{stewartred}8.}\ \, $\displaystyle \lim_{x \to a} c = c$ \hspace{0.30\textwidth} \textbf{\color{stewartred}9.}\ \, $\displaystyle \lim_{x \to a} x = a$
    \end{tcolorbox}

    Các giới hạn này là hiển nhiên theo quan điểm trực quan (hãy diễn đạt chúng bằng lời hoặc vẽ đồ thị của $y = c$ và $y = x$), nhưng các phép chứng minh dựa trên định nghĩa chính xác được yêu cầu trong các Bài tập 2.4.23--24.

    Bây giờ nếu đặt $f(x) = x$ trong Quy tắc 6 và sử dụng Quy tắc 9, ta nhận được một giới hạn đặc biệt hữu ích cho các hàm lũy thừa.

    % Khung 10
    \begin{tcolorbox}[
        sharp corners,
        colback=white,
        colframe=stewartred,
        boxrule=0.75pt,
        left=10pt, right=10pt, top=5pt, bottom=5pt,
        before skip=6pt, after skip=6pt
    ]
    \textbf{\color{stewartred}10.}\ \, $\displaystyle \lim_{x \to a} x^n = a^n$ \qquad\quad trong đó $n$ là một số nguyên dương
    \end{tcolorbox}
\end{minipage}
